\section{Examples and the contract}
\label{sec:lvexamples}

All fits in this section use one protocol --- the grids of
\cref{sec:lvinfo}, the objective of \cref{sec:lvcalib}, every constant
recorded in \cref{app:lvnumerics} --- applied unchanged to a synthetic
market whose answer is known and to the book's frozen SPY and NVDA
chains.  No fit is tuned separately.

\subsection{The round trip, and its two separate numbers}

The synthetic market prices a known smooth field through an
\emph{independent} dense scheme (so the exercise cannot flatter the
chapter's march by validating it against itself), quotes it at four
expiries, and hands the \MacLvRtNQuotes\ clean quotes to the calibration
from a flat seed.  Per \cref{sec:lvinfo}, the result is two numbers, not
one.  The \emph{quotes} reprice to \MacLvRtRmsErrBp\,bp rms and
\MacLvRtMaxErrBp\,bp at worst --- the fit is essentially exact in the
space the optimizer sees.  Separately, the \emph{field itself} agrees with
the truth to \MacLvRtSurfRmsPts\ volatility points rms on the
quote-covered region, with a maximum of \MacLvRtSurfMaxPts\ points ---
more than a hundred times the quote error, concentrated exactly where
\cref{fig:lvidentify} said the quotes do not look: between expiries and
near the span edges (\cref{fig:lvrecovery}b).  Both numbers are honest;
only reporting the first would not be.

\begin{figure}[htbp]
\centering
\paperfig{\textwidth}{fig_lv_recovery.pdf}
\caption{The synthetic round trip.  (a)~Target quotes (dots) and the
recovered fit (lines) at the four expiries: rms \MacLvRtRmsErrBp\,bp, max
\MacLvRtMaxErrBp\,bp.  (b)~The second number: $|$fitted $-$ true$|$ local
volatility over the quote-covered region, quote positions dotted.  The
error lives between the expiries and at the edges of the quoted spans ---
largest where the nearest quote is farthest.}
\label{fig:lvrecovery}
\end{figure}

\subsection{The frozen chains}

The real test fits one sheet per underlier to every prepared mid quote of
the snapshot's eight expiries at once: \MacLvSpyQuotes\ quotes against
\MacLvSpyVtx\ nodal variances for SPY, \MacLvNvdaQuotes\ against
\MacLvNvdaVtx\ for NVDA.  \Cref{fig:lvsheet} already showed the SPY sheet
itself; \cref{fig:lvfit} shows four of its expiries against the quotes,
priced by the \emph{one} calibrated diffusion --- a two-day smile and a
year-and-a-half LEAP from the same list of positive numbers.

\begin{figure}[htbp]
\centering
\paperfig{\textwidth}{fig_lv_fit.pdf}
\caption{Four SPY expiries priced by the single calibrated sheet, against
the frozen quotes (dots: mids; whiskers: bid--ask).  Panel titles give the
expiry and its distance; annotations the per-expiry rms on the fitting
operator.}
\label{fig:lvfit}
\end{figure}

\begin{figure}[htbp]
\centering
\paperfig{\textwidth}{fig_lv_rms.pdf}
\caption{Per-expiry rms for SPY (a) and NVDA (b), on the fitting operator
(grey) and on the refined operator (blue), log scale.  The refined bars
are the ones to read: the gap between the pairs is the compensated
discretization error of \cref{sec:lvlattice}, and it owns the short end.}
\label{fig:lvrms}
\end{figure}

\Cref{fig:lvrms} reports the chapter's two metrics side by side.  On the
fitting operator the surfaces read \MacLvSpyRmsBp\,bp (SPY) and
\MacLvNvdaRmsBp\,bp (NVDA) rms across all quotes --- NVDA's wider spreads
and harder short-dated smiles carry more residual everywhere.  Repriced on
the refined operator the same surfaces read \MacLvSpyConvBp\ and
\MacLvNvdaConvBp\,bp, and the worst single expiry --- in both names the
one-day front --- reads \MacLvSpyWorstConvBp\ and
\MacLvNvdaWorstConvBp\,bp.  The gap between the paired bars is the point,
and it is the operator's blind spot of \cref{sec:lvlattice} made
quantitative: the optimizer partially cancels the coarse march's error at
the expiries it is fitted on, the in-operator residual flatters
accordingly, and the number to quote is the refined one.  Where the two
agree (the long end) the fit is genuinely converged; where they diverge
(day-scale maturities) the remaining error is discretization, not
disagreement with the quotes.  The measured cleanliness of
\cref{rem:lvscope} completes the report: across every expiry of both
calibrated sheets, the minimum divided second difference of the marched
prices is \MacLvSpyButterflyMin\ (SPY) and \MacLvNvdaButterflyMin\ (NVDA),
and the minimum adjacent-expiry increment is \MacLvSpyCalendarMin\ and
\MacLvNvdaCalendarMin\ --- solver rounding, on both counts, as
\cref{prop:lvpropagate} predicts and \cref{rem:lvscope} declines to assume.

\subsection{The contract}
\label{sec:lvcontract}

\begin{table}[htbp]
\centering\small
\begin{tabular}{@{}>{\raggedright\arraybackslash}p{0.31\textwidth}
>{\raggedright\arraybackslash}p{0.31\textwidth}
>{\raggedright\arraybackslash}p{0.30\textwidth}@{}}
\toprule
\textbf{Structural} & \textbf{Numerically controlled} &
\textbf{Outside the claim}\\
\midrule
Box bounds imply a positive field on the hull; a positive field implies
butterfly-free, calendar-ordered marginals (\cref{prop:lvnodal,%
prop:lvpropagate}); every optimizer iterate is an admissible diffusion &
Butterfly, calendar and bound cleanliness of the marched prices, measured
per fit; refined-operator reprice beside the in-operator residual &
Convexity preservation of the discrete scheme as a theorem; the
extrapolation region beyond the hull, where basis weights are signed\\[4pt]
Monotone, unconditionally $\ell^\infty$-stable march
(\cref{prop:lvmmatrix,cor:lvdmp}); exact tangent Jacobian from the same
factorization (\cref{prop:lvtangent}) & Tangent columns audited against
finite differences; an independent discretization cross-checks the march &
Uniqueness of the calibrated sheet: the data determines an equivalence
class, the conventions of \cref{sec:lvcalib} a representative\\[4pt]
Bounded field $\Rightarrow$ implied variance bounded in strike, strictly
inside Lee's cap on the hull (\cref{rem:lvlee}) & Grid placement rules
resolving each expiry's $\sigma\sqrt\tau$; adaptive variance box keyed to
ATM quotes & Jumps and scheduled-event spikes (the clock, Chapter~8);
spot--volatility dynamics (Chapter~9); the deep-put tail's shape without a
variance-swap quote (Chapter~5)\\
\bottomrule
\end{tabular}
\caption{The local-volatility model contract: theorem, numerical control,
and excluded scope.}
\label{tab:lvcontract}
\end{table}

The chapter's separation thesis reads off the table.  \emph{Validity} is
structural and cheap: a box.  \emph{Information} is scarcer than the
formula \eqref{eq:lvextract} pretends: quotes determine the field only
through smoothing integrals, and only where they exist.  \emph{Convention}
fills the rest --- roughness, reference, tie, wing, seed.  The calibration
does not eliminate those choices; it declares them, measures what the
discretization leaks, and reports the error twice: once as the optimizer
saw it, once as it is.
